\subsection{Continuous Variables and Constraints}\label{sec:continuous-variables-and-constraints}

\subsubsection{Continuous Variables through Binary Expansion}\label{sec:continuous-variables-through-binary-expansion}

This section addresses the first practical limitation of QUBO: \textbf{variables must be binary}. Up to now this was not a problem because we were working with Boolean variables $a,b \in \{0,1\}$. However, most real-world optimization problems involve quantities such as temperatures, distances, costs, production levels, or resource allocations, which are usually \textbf{represented by continuous or integer-valued variables}. The natural question is therefore, ``\emph{how can a QUBO formulation represent variables that are not binary?}''. The answer is \textbf{binary expansion}.

\highspace
\begin{flushleft}
    \textcolor{Green3}{\faIcon{book} \textbf{Binary Expansion}}
\end{flushleft}
Recall how ordinary binary numbers are represented. For example, the number $13$ can be represented in binary as:
\begin{equation*}
    13 = 1 \cdot 2^3 + 1 \cdot 2^2 + 0 \cdot 2^1 + 1 \cdot 2^0 = \left(1101\right)_{2}
\end{equation*}
\hl{The same principle can be used inside a QUBO formulation.} Instead of representing a variable $x$ directly, we write:
\begin{equation}\label{eq:binary-expansion}
    x =\sum_{k} 2^k x_k \qquad x_{k} \in \{0,1\}
\end{equation}
Each binary variable $x_k$ determines whether the corresponding power of two contributes to the value of $x$.

\highspace
Suppose we want to represent a real variable using binary variables. One possible encoding is:
\begin{equation*}
    x = \sum_{k=-2}^{2} 2^{k} x_{k}
\end{equation*}
Where $k$ ranges from $-2$ to $2$. Expanding the powers of two, we get:
\begin{equation*}
    x = \frac{1}{4} x_{-2} + \frac{1}{2} x_{-1} + x_{0} + 2 x_{1} + 4 x_{2}
\end{equation*}
Each coefficient corresponds to a different binary weight.
\begin{itemize}
    \item \textbf{\emph{What values can be represented?}} Since each variable is binary $x_i \in \{0,1\}$, the smallest representable value is obtained when all variables are zerp, $x_{\min} = 0$. The largest representable value is obtained when all variables are one, $x_{k} = 1$ for all $k$, which gives:
    \begin{equation*}
        x_{\max} = \frac{1}{4} + \frac{1}{2} + 1 + 2 + 4 = \frac{31}{4} = 7.75
    \end{equation*}
    Therefore the representation can descibe values between $0$ and $7.75$.

    \item \textbf{\emph{Precision}}. The smallest coefficient is $\frac{1}{4}$, as a consequence, the representable values differ by multiplies of $0,25$. This quantity is called the \textbf{resolution} (or \emph{precision}) of the representation. Adding more negative powers of two would increase the precision, while adding more positive powers of two would increase the maximum representable value.
\end{itemize}
\textcolor{Green3}{\faIcon{question-circle} \textbf{Who decides the range and precision of the representation?}} The \textbf{modeler decides both the expansion and the precision}. For instance, suppose our original optimization problem contains a variable $x \in \left[0,5\right]$. To convert the problem into QUBO, we must decide how to encode $x$ using binary variables. If we choose the same expansion as before, we can represent values up to $7.75$, which is sufficient to cover the range of $x$. In general, \textbf{higher precision requires more binary variables}, which \textbf{increases the size of the QUBO problem} and can make it \textbf{more difficult to solve}. Therefore, there is a trade-off between precision and computational complexity that the modeler must consider when choosing the binary expansion.

\highspace
In general, \definition{Binary Expansion} is simply the idea of \textbf{representing a number as a sum of powers of two}. It is the binary equivalent of decimal notation. In the context of QUBO, it allows us to \textbf{replace one continuous} (or integer) \textbf{variable with several binary variables}.